% 09-12-discusion-conclusiones.tex — Fuente: docs/informe-final/09-12-discusion-conclusiones.md

\chapter{Discusión}
\label{cap:discusion}

\section{Respuesta a las preguntas de investigación}

\textbf{RQ1 --- ¿Cumple el sistema los umbrales de rendimiento, seguridad,
usabilidad, cobertura y calidad web declarados?} Parcialmente. De los
cinco bloques evaluados (Capítulo~\ref{cap:evaluacion-resultados}), cuatro
cumplen sus umbrales por completo (rendimiento, seguridad
manual+automatizada, usabilidad, calidad web en ambos perfiles). El
bloque de cobertura cumple en líneas (72.0\,\% $\geq$70\,\%) pero
\textbf{no en ramas} (62.5\,\% $<$70\,\%) --- una respuesta honesta a RQ1
no puede ser ``sí'' sin matiz. El sistema cumple la gran mayoría de sus
propiedades de calidad declaradas con evidencia reproducible, con una
brecha cuantificada y no oculta en cobertura de ramas.

\textbf{RQ2 --- ¿Es la estrategia híbrida de acceso a datos implementable
y verificable dentro de las restricciones del proyecto?} Sí, parcialmente
y con matices declarados. Se conectaron 7 procedimientos/funciones
end-to-end al código en ejecución (por encima del mínimo de 6 exigido por
la guía), sin concatenación de SQL dinámico detectada por análisis
estático ni por auditoría manual (Capítulo~\ref{cap:evaluacion-resultados}~\S8.2).
Sin embargo, las 6 rutinas originales de la Tercera Entrega siguen sin
conectar (Capítulo~\ref{cap:diseno-arquitectura}, ADR-006), y existe una
nuance de conformidad literal sobre el mecanismo de invocación
(\texttt{@Query(nativeQuery=true)} parametrizado en vez de
\texttt{@Procedure}/\texttt{@NamedStoredProcedureQuery} estrictos,
Capítulo~\ref{cap:implementacion}~\S7.4) que el equipo debe resolver antes
del cierre.

\textbf{RQ3 --- ¿Produce el proceso de ingeniería de requisitos aplicado
un corpus verificable y trazable?} Sí, con una brecha de conformidad
sintáctica declarada. El 100\,\% de los 37 requisitos tiene trazabilidad
completa hacia código y pruebas (\texttt{matriz.csv}, validado en CI), y
el 89.2\,\% está implementado o verificado. Nueve de las quince
características de calidad INCOSE~v4 se cumplen plenamente, pero el
corpus no sigue el patrón sintáctico formal de
ISO/IEC/IEEE~29148:2018 (característica C9), y dos requisitos violan
singularidad (C5) --- hallazgos concretos, no una declaración genérica de
``cumple parcialmente'' (Capítulo~\ref{cap:ingenieria-requisitos}~\S4.3).

\textbf{RQ4 --- ¿Qué amenazas a la validez afectan la generalización de
los resultados?} Ver Capítulo~\ref{cap:validez} para el análisis completo
por categoría; en resumen, la amenaza más relevante es de validez interna
en la comparación caché frío/caliente (diseño experimental que no aísla
la variable de interés) y de validez externa en el muestreo por
conveniencia del bloque de usabilidad.

\section{Comparación con trabajos relacionados}

Dado que el Capítulo~\ref{cap:trabajos-relacionados} declara honestamente
la ausencia de una revisión sistemática de literatura completa, esta
comparación se limita a lo que el capítulo sí pudo establecer: ninguna de
las referencias contextuales citadas ---incluidas las plataformas
multilaterales~\cite{HAGIU2015}, los protocolos de \emph{escrow}
verificable~\cite{ASGAONKAR2019} y los estudios de microservicios en
producción~\cite{WANG2021}--- combina patrón \emph{escrow}, verificación
de identidad asistida por IA y estrategia híbrida de acceso a datos
evaluada empíricamente en un mismo sistema (\S3.3). Una comparación
cuantitativa rigurosa (tabla comparativa de $\geq$8 filas con métricas
equivalentes) requiere primero completar la revisión de literatura
pendiente --- no se fabrica aquí una comparación que la evidencia
disponible no sostiene.

\section{Hallazgos inesperados}

El hallazgo más relevante no anticipado al inicio del proyecto fue el
propio del bloque de rendimiento (\S8.1): el escenario ``frío'' resultó
con mejor latencia que el ``caliente'', revelando un defecto de diseño del
script de carga (no del sistema evaluado) que invalida la comparación tal
como está construida. Detectarlo y documentarlo con honestidad, en vez de
descartarlo o presentarlo como validado, es en sí mismo parte del aporte
metodológico de este trabajo.

\section{Implicaciones prácticas}

Para equipos que evalúen adoptar una estrategia híbrida de acceso a datos
similar: la evidencia de este proyecto sugiere que conectar procedimientos
almacenados end-to-end (no solo declararlos en SQL) consume un esfuerzo
de verificación no trivial --la brecha entre ``el archivo \texttt{.sql}
existe'' y ``el código Java realmente lo invoca'' fue precisamente el
hallazgo que motivó la ampliación de \texttt{ADR-006} documentada en el
Capítulo~\ref{cap:diseno-arquitectura}-- y que declarar la nuance de
conformidad del mecanismo de invocación
(Capítulo~\ref{cap:implementacion}~\S7.4) antes de una auditoría externa
evita sorpresas en la evaluación.

\chapter{Amenazas a la validez}
\label{cap:validez}

\section{Validez de constructo}

¿Miden los instrumentos elegidos lo que pretenden medir? El instrumento
SUS mide usabilidad percibida, no usabilidad objetiva (tiempo de tarea,
tasa de error) --- el proyecto no midió usabilidad objetiva como
complemento, por lo que la conclusión de ``usabilidad aceptable'' se
apoya en un único constructo subjetivo. \textbf{Mitigación aplicada:} se
usó el instrumento SUS sin modificar (10 ítems originales de
Brooke~\cite{BROOKE1996}), evitando la amenaza de validez de constructo
que introduciría una versión ad-hoc del cuestionario.

\section{Validez interna}

La comparación caché frío/caliente del bloque de rendimiento (\S8.1,
\S9.1) tiene una amenaza de validez interna real y ya documentada: el
diseño del script de carga no varía los parámetros de consulta entre
iteraciones, por lo que la variable independiente (estado del caché) no
se manipula de forma efectiva durante la mayoría de las 1500 iteraciones
por corrida. \textbf{Mitigación aplicada:} declarar la limitación
explícitamente en vez de presentar la comparación como válida
(\texttt{docs/mediciones/perf/REPORTE-PERF.md}); \textbf{mitigación
pendiente:} rediseñar el script para variar parámetros de forma
controlada y forzar \emph{cache misses} reales y medibles.

\section{Validez externa}

El bloque de usabilidad usa una muestra de conveniencia (N=16, red de
contactos del equipo, ver Capítulo~\ref{cap:materiales-metodos}~\S5.4), no
una muestra aleatoria de la población objetivo real (creadores y clientes
de economía creativa). Siguiendo la orientación de Baltes y
Ralph~\cite{BALTES2022} sobre reporte de muestreo en ingeniería de
software empírica, esta limitación se declara explícitamente en vez de
generalizar el resultado SUS a ``los usuarios de Artisync'' sin matiz: el
resultado (76.88, categoría~B de Bangor~\cite{BANGOR2009}) es válido para
la muestra evaluada, con generalización externa limitada a un perfil
demográfico similar (adultos con experiencia web media-alta, dispositivos
variados).

\section{Validez de conclusión}

Las conclusiones de los bloques de rendimiento y cobertura se apoyan en
estadística descriptiva (media, DT, IC~95\,\%) sin test inferencial,
porque no se identificaron comparaciones de dos condiciones
estadísticamente válidas para probar dentro del alcance actual (la única
candidata, caché frío/caliente, está invalidada por la amenaza de validez
interna de \S10.2). Esto es consistente con la recomendación de preferir
IC~95\,\% sobre valores p como reporte primario
(Capítulo~\ref{cap:materiales-metodos}~\S5.6), siguiendo el protocolo
estadístico de Arcuri y Briand~\cite{ARCURI2011} para ingeniería de
software, pero implica que ninguna afirmación de ``diferencia
significativa'' se sostiene en este documento --- y en efecto, ninguna se
hace.

\chapter{Trabajo futuro}
\label{cap:trabajo-futuro}

\begin{enumerate}
  \item \textbf{Conectar las 6 rutinas SQL originales de la Tercera
    Entrega} (\texttt{fn\_catalogo\_filtrado},
    \texttt{fn\_calificacion\_promedio\_creador},
    \texttt{fn\_reporte\_comisiones\_creador},
    \texttt{fn\_cerrar\_pedidos\_vencidos},
    \texttt{fn\_liberar\_fondos\_escrow},
    \texttt{fn\_generar\_codigo\_pedido}) al código Java en ejecución,
    cerrando la deuda técnica admitida en \texttt{ADR-006}. Quedó fuera
    del alcance de la ampliación del 16 de agosto porque esa ampliación
    priorizó sumar rutinas nuevas del módulo de seguridad, no reparar las
    existentes.
  \item \textbf{Rediseñar el script de carga k6} para forzar \emph{cache
    misses} reales y medibles por iteración, habilitando un test
    inferencial (Wilcoxon) con tamaño de efecto (\emph{Cliff's delta})
    válido sobre la comparación caché frío/caliente~\cite{ARCURI2011} ---
    hoy inválida por la amenaza de validez interna documentada en
    \S10.2.
  \item \textbf{Elevar la cobertura de ramas de JaCoCo} del 62.5\,\%
    actual al umbral de 70\,\% exigido, con desglose por módulo
    (dominio/servicio/controlador) en vez del agregado global actual, y
    activar \texttt{jacoco:check} en \texttt{pom.xml} para que el build
    falle bajo el umbral en vez de solo reportarlo.
  \item \textbf{Ejecutar una revisión de literatura reducida} (2--3 bases
    de datos, 10--15 términos clave) para completar el
    Capítulo~\ref{cap:trabajos-relacionados} con la tabla comparativa de
    $\geq$8 filas que la guía exige y que este documento declaró
    honestamente como pendiente.
  \item \textbf{Resolver la nuance de conformidad de invocación de
    procedimientos} (Capítulo~\ref{cap:implementacion}~\S7.4): decidir
    entre reescribir las 7 funciones nuevas como \texttt{PROCEDURE} para
    usar \texttt{@Procedure} en sentido estricto, adoptar
    \texttt{@NamedStoredProcedureQuery}, o documentar formalmente por qué
    \texttt{@Query(nativeQuery=true)} parametrizado satisface el
    requisito.
\end{enumerate}

\chapter{Conclusiones}
\label{cap:conclusiones}

Este trabajo diseñó, implementó y evaluó empíricamente Artisync, una
plataforma web de intermediación entre creadores de contenido digital y
clientes, con una estrategia híbrida de acceso a datos y verificación de
identidad asistida por inteligencia artificial. Respondiendo brevemente a
las preguntas de investigación planteadas en el
Capítulo~\ref{cap:introduccion}: el sistema cumple la mayoría de sus
umbrales de calidad declarados con evidencia empírica reproducible (RQ1),
con una brecha cuantificada en cobertura de ramas; la estrategia híbrida
de acceso a datos es implementable y verificable dentro de las
restricciones de un proyecto académico de 17 semanas, con 7
procedimientos conectados end-to-end y una nuance de conformidad
declarada sobre el mecanismo de invocación (RQ2); el proceso de
ingeniería de requisitos produjo un corpus 100\,\% trazable hacia código
y pruebas, con brechas concretas de conformidad sintáctica INCOSE~v4
(RQ3); y las amenazas a la validez más relevantes ---diseño del
experimento de caché y muestreo de usabilidad por conveniencia--- están
documentadas y mitigadas en la medida de lo posible dentro del alcance
del proyecto (RQ4).

Las cinco contribuciones enumeradas en el
Capítulo~\ref{cap:introduccion}~\S1.4 se sostienen con evidencia concreta
de los Capítulos~\ref{cap:diseno-arquitectura} a
\ref{cap:evaluacion-resultados}: el sistema software completo y
desplegable (Capítulos~\ref{cap:diseno-arquitectura}--\ref{cap:implementacion}),
la estrategia híbrida de acceso a datos verificada end-to-end
(Capítulo~\ref{cap:implementacion}~\S7.4), el corpus de ingeniería de
requisitos trazable (Capítulo~\ref{cap:ingenieria-requisitos}), el
paquete de evidencia empírica reproducible
(Capítulo~\ref{cap:evaluacion-resultados}, con procedencia declarada en
\texttt{docs/mediciones/DATA-PROVENANCE.md}), y los siete ADR
(Capítulo~\ref{cap:diseno-arquitectura}~\S6.2).

El valor distintivo de este documento, más allá de reportar únicamente
resultados favorables, es la disciplina de declarar honestamente cada
brecha encontrada durante su propia redacción --desde la cobertura de
ramas por debajo del umbral hasta la nuance de conformidad de invocación
de procedimientos-- siguiendo el mismo estándar de auditoría verificable
que produjo \texttt{docs/observaciones/INFORME-BRECHAS-ENTREGA-FINAL.md} y
que este documento hereda explícitamente.
